%% kapa: exam=1 round=37 year=2026 period=2 form=A subject=appraisal_law subject_ko=감정평가관계법규 seq=16 number=16
%% source: Q-Net 감정평가사 기출문제 (cst00309) artlSeq=5258924 file=raw/1008/5258924/2026년 제37회 감정평가사 1차 2교시.pdf
%% conversion: kordoc 4.24.51 -> tools/gen1.py (auto); review=proofread-sonnet; flags=none
\begin{question}{16}
부동산 가격공시에 관한 법령상 개별공시지가에 관한 설명으로 옳지 않은 것은?
\begin{choices}
  \item 개별공시지가 조사ㆍ산정의 기준에는 토지가격비준표의 사용에 관한 사항이 포함되어야 한다.
  \item 개별공시지가를 결정ㆍ공시하기 위하여 개별토지가격 산정의 타당성 검증을 의뢰받은 감정평가법인은 비교표준지 선정의 적정성에 관한 사항을 검토ㆍ확인하고 의견을 제시해야 한다.
  \item 농지보전부담금의 부과대상이 아닌 농지에 대하여는 개별공시지가를 결정ㆍ공시하여야 한다.
  \item 시장ㆍ군수 또는 구청장은 개별공시지가에 대한 이의신청을 심사하기 위하여 필요할 때에는 감정평가법인등에게 검증을 의뢰할 수 있다.
  \item 관계 공무원은 일출 전ㆍ일몰 후에는 토지 점유자의 승인 없이 개별공시지가의 산정을 위하여 울타리로 둘러싸인 타인의 토지에 출입할 수 없다.
\end{choices}
\end{question}
\answer{3}
